\section*{How to read this paper}
\addcontentsline{toc}{section}{How to read this paper}
\markboth{How to read this paper}{}

This paper is written for two readers at once. One has never studied algebra beyond
school and wants to understand what has actually been found. The other wants every
statement to be precise enough to check. We have tried not to cheat either of them.
The mathematics is stated in full; next to it, in gold-edged boxes, the same idea is
explained in ordinary words, and in teal-edged boxes we say what the idea means
physically. Rose-edged boxes mark the edge of what is known.

\subsection*{The five badges}

Every claim of substance carries one of five labels. They are the most important
typographical device in the paper, because the whole value of a speculative subject
lies in being able to tell the proved from the hoped-for.

\begin{center}
\renewcommand{\arraystretch}{1.35}
\begin{tabularx}{0.95\textwidth}{@{}lX@{}}
\toprule
\tierT & A mathematical theorem: proved by an argument in the text, or a standard result
         with a reference. \\
\tierC & An exact finite computation, carried out by a committed script whose output
         certificate is named. ``Exact'' means integer or rational arithmetic, or an
         exhaustive enumeration --- never a numerical approximation, unless we say so. \\
\tierB & A bridge: a mathematically true statement \emph{plus} an interpretation of it
         (for example, calling a mathematical light cone a light cone of time). The
         mathematics is certain; the interpretation is a proposal. \\
\tierO & An open question, stated as sharply as we can. \\
\tierX & Something that was once believed, written down, and then shown false. We keep
         these, with the reason. \\
\bottomrule
\end{tabularx}
\end{center}

\subsection*{A map of the journey}

The paper follows one chain of exact identifications. Figure~\ref{fig:journey} is the
whole argument in one picture; each box is a section.

\begin{figure}[htbp]
\centering
\begin{tikzpicture}[
  node distance=5.5mm and 16mm,
  st/.style={draw=ink!60,rounded corners=3pt,fill=white,align=center,
    font=\sffamily\footnotesize,text width=50mm,minimum height=10mm,inner sep=3pt},
  hl/.style={st,fill=tealsoft,draw=teal},
  ar/.style={-{Stealth[length=2.4mm]},thick,ink!55}]
  \node[st] (q) {one qutrit $\C^3$\\\scriptsize three settings, with phases};
  \node[st,below=of q] (h) {9 operator histories $\End(\C^3)$\\\scriptsize the qutrit acting on itself};
  \node[hl,below=of h] (w) {$\W$: 40 points, 40 lines\\\scriptsize \S2--\S3};
  \node[st,below=of w] (m) {27-event finite space-time\\\scriptsize light cone $t^2-x^2-y^2$, \S4};
  \node[st,below=of m] (c) {four clocks $=\mathbb{P}^1(\F_3)$\\\scriptsize tetracode, Hesse striations, \S4};
  \node[hl,right=of q] (e8) {$E_8$: 248 dimensions\\\scriptsize $80+84+84$ from 9 histories, \S5};
  \node[st,below=of e8] (gr) {clock gradings of $E_8$\\\scriptsize $1{+}56{+}134{+}56{+}1$; \ cubic ladder $54\to81\to83$};
  \node[st,below=of gr] (f) {Freudenthal $56$ and quartic\\\scriptsize the $57$-dimensional light cone, \S6};
  \node[st,below=of f] (r) {real forms: $E_{8(-24)}$ vs $E_{8(8)}$\\\scriptsize where signature comes from, \S7};
  \node[hl,below=of r] (l) {Albert algebra $\to\mathfrak{so}(1,9)$\\\scriptsize Weyl spinor 16 $=$ matter, \S8};
  \draw[ar] (q)--(h); \draw[ar] (h)--(w); \draw[ar] (w)--(m); \draw[ar] (m)--(c);
  \draw[ar] (c.east) -- ++(8mm,0) |- (e8.west);
  \draw[ar] (e8)--(gr); \draw[ar] (gr)--(f); \draw[ar] (f)--(r); \draw[ar] (r)--(l);
\end{tikzpicture}
\caption{The chain of exact identifications followed in this paper. Shaded boxes are the
three places where the story changes scale: from a qubit-sized object to a finite
geometry, from the geometry to $E_8$, and from $E_8$ to a Lorentzian algebra.}
\label{fig:journey}
\end{figure}

\subsection*{What you will need}

Nothing beyond patience. We introduce every structure before we use it, and the
glossary in Appendix~B collects the vocabulary. Readers who want only the story can
read the plain-words boxes and the key-idea panels and skip the displayed mathematics;
nothing essential will be lost. Readers who want only the mathematics can do the
reverse.

\subsection*{Where the material comes from}

Everything here is drawn from the public repository \texttt{wilcompute/W33-Theory},
a record of several thousand numbered ``passes'' by two independent research tracks,
and from its three long manuscripts --- the $\W$ paper, the photonic Holonet paper and
the Holonet machine blueprint --- some of whose pages are adapted here. Passes 10949 and
10950, which supply Sections 6--8, were carried out for this synthesis. A reader who
wants to check a statement can find its script and certificate in the ledger of
Appendix~A.
